\chapter*{术语索引}
\phantomsection
\addcontentsline{toc}{chapter}{术语索引}

参考文献号依次表示章、节与小节（个别例外情形为习题）。

测度的绝对值：III, 1, 6 与 VI, 2, 8。

适应对，\(\mu-\)：V, 4, 1。

加性集函数：IV, 4, 9。

加性，完全：IV, 4, 5。

适宜映射，\(\mu-\)：V, 3, 1 与 VI, 1, 1，脚注。

互奇异测度：V, 5, 7。

几乎处处，有定义的函数：IV, 2, 5。

几乎处处，成立的性质：IV, 2, 3。

原子测度：III, 1, 3。

带，完全格序空间中的：II, 1, 5。

测度的重心：IV, 7, 1。

以 \(\mu\) 为基的测度：V, 5, 2 与 VI, 2, 8。

标量意义以 \(\mu\) 为基（向量测度）：VI, 2, 5。

以 \(\mu\) 为基的向量测度：VI, 2, 4。

属于扩散定义域的测度：V, 3, 5。

Bishop 定理：IV, 7, 5。

边界（边缘）：III, 1, 1，脚注。

有界扩散：V, 3, 5。

依测度有界的函数：IV, 6, 2 与 IV, 6, 3。

有界测度：III, 1, 8 与 VI, 2, 9。

承载测度的子集：V, 5, 7。

Choquet 定理：IV, 7, 2 与 IV, 7, 6。

一个集的子集组成的集环（集族）：IV, 4, 9。

等价类（函数的），关于一测度的：IV, 2, 4、IV, 2, 5 与 IV, 5, 2。

伪像类（测度类的伪像）：VI, 3, 2。

余格子空间：II, 1，习题 3。

紧收敛（在测度空间中）：III, 1, 10。

完全加性：IV, 4, 5。

复测度：III, 1, 3 与 VI, 2, 8。

复合扩散：V, 3, 6。

集中于一集上的测度：V, 5, 7。

共轭指数：IV, 6, 4。

共轭测度：III, 1, 5 与 VI, 2, 8。

连通分支：V, 6，习题 12，脚注。

保守映射：V, 7，习题 9。

几乎处处连续：IV, 5，习题 16。

平均收敛、\(p\) 阶平均收敛、均方收敛：IV, 3, 3。

依测度收敛：IV, 5, 11。

收敛，几乎处处：IV, 2, 5。

收敛，紧（在测度空间中）：III, 1, 10。

收敛，严格紧（在测度空间中）：III, 1, 10。

收敛，淡：III, 1, 9。

完全格序空间中的收敛滤子：II, 1，习题 9。

无穷远处可数：III, 1, 1，脚注。

可数凸性定理：IV, 3, 2。

测度的切片分解：V, 6, 6 与 VI, 3. 1。

函数的递减重排：IV, 5，习题 29。

几乎处处有定义的（函数）：IV, 2, 5。

局部几乎处处有定义的（函数）：IV, 5, 2。

以密度 \(g\) （关于 \(\theta\) ）定义的测度：V, 5, 2 与 VI, 2, 8。

密度点（子集的）：V, 6，习题 15。

密度，向量测度关于正测度的：VI, 2, 4。

密度，关于复测度的：III, 1, 4、V, 5, 2 与 VI, 2, 8。

函数的 Dini 导数：V, 6，习题 12。

弥散测度：V, 5, 10。

扩散：V, 3, 5。

扩散，有界：V, 3, 5。

扩散，复合：V, 3, 6。

Dirac 测度：III, 1, 3。

直接极限，局部凸空间直接系统的：III, 1, 1。

直和，序：II, 1, 4。

离散测度：III, 1, 3。

测度的分解（测度分解）：V, 6, 6 与 VI, 3, 1。

分解（测度分解），测度 \(\mu\) 关于 \(\mu\) -逆紧映射的：VI, 3, 1。

分解（测度分解），测度 \(\mu\) 关于 \(\mu\) 的伪像的：VI, 3, 3。

分解（测度分解），按可测等价关系对测度作的：VI, 3, 5。

分布，序列的极限：III, 2，习题 5。

定义域，扩散的，属于其定义域的测度：V, 3, 5。

二重积分：III, 4, 1。

Dunford--Pettis 定理：VI, 2, 5。

Egoroff 定理：IV, 5, 4。

等度可积集（函数的）：IV, 5, 11。

等分布序列（关于一测度的）：III, 2，习题 5。

等可测函数：IV, 5，习题 29。

关于一测度的函数等价类：IV, 2, 4、IV, 2, 5 与 IV, 5, 2。

等价关系，Hausdorff：VI, 3, 4。

等价关系，可测：VI, 3, 4。

等价测度：V, 5, 6 与 VI, 2, 8。

等价，\(\mu\) -等价（函数）：IV, 2, 4。

本质上积分：V, 1, 1。

本质 \(\mu\) -可积（集、函数）：V, 1, 3。

本质可积函数：V, 1, 3。

本质可积函数，p 次幂：V, 1, 3。

本质可积函数，关于向量测度的：VI, 2, 2。

在 A 上本质可积的函数：V, 5, 3。

膨胀集（右、左），函数的：V, 6，习题 12。

极值点，\(\mathcal{H}-\)：IV, 7, 3。

几乎处处有限的函数：IV, 2, 6。

线性形式，正：II, 2, 1。

线性形式，相对有界：II, 2, 2。

完全格序空间：II, 1, 3。

只依赖于有限个变量的函数：III, 4, 6。

函数，在 A 上 \(\mu\) -可测的：V, 5, 3。

函数，\(\mu\) -适度的：V, 1, 2。

函数，\(\Phi\) -阶梯：IV, 4, 9。

函数，\(p\) 次幂本质可积的：V, 1, 3。

函数，\(p\) 次幂可积的：IV, 3, 4。

函数，依测度有界的：IV, 6, 2 与 IV, 6, 3。

函数，几乎处处有定义的：IV, 2, 5。

函数，局部几乎处处有定义的：IV, 5, 2。

函数，本质可积：V, 1, 3。

函数，关于向量测度本质可积的：VI, 2, 2。

函数，在 A 上本质可积的：V, 5, 3。

函数，几乎处处有限的：IV, 2, 6。

函数，标量意义具有某性质的：VI, 1, 1。

函数，可积的，\(\mu\) -可积的：IV, 4, 1。

函数，Lebesgue 第 n 个：IV, 6，习题 15。

函数，局部可积：V, 5, 1。

函数，在 A 上局部可积的：V, 5, 3。

函数，局部可忽略的，局部 \(\mu\) -可忽略的：IV, 5, 2。

函数，可测阶梯：IV, 5, 5。

函数，可测的，\(\mu\) -可测的：IV, 5, 1。

函数，可测的，定义在 X 的可测子集上的：IV, 5, 10 与 V, 5, 3。

函数，可忽略的，\(\mu\) -可忽略的：IV, 2, 1 与 IV, 2, 4。

函数，数值：I, 1, 1，脚注。

函数，标量本质可积：V, 3, 1、VI, 1, 1 与 VI, 2, 10。

函数，标量良可积：VI, 1，习题 19。

函数，阶梯的，\(\Phi\) -阶梯：IV, 4, 9。

函数，支集：III, 1, 1。

函数，普遍可测：V, 3, 4。

函数，淡 \(\mu\) -可测：V, 3, 1。

函数，淡连续：V, 3, 1。

函数，等度可积的：IV, 5, 11。

函数，等可测的：IV, 5，习题 29。

函数，等价的，\(\mu\) -等价的：IV, 2, 4。

性质（GDF）：VI, 1, 4。

Gelfand--Dunford 定理：VI, 1, 4。

滑峰法：V, 5，习题 13。

\(\mathcal{H}\) -极值点：IV, 7, 3。

Hölder 不等式：I, 2 与 IV, 6, 4。

Haar 标准正交系：IV, 6，习题 17。

Hardy 不等式：IV, 6，习题 19。

Hausdorff 等价关系：VI, 3, 4。

测度的像：V, 6, 1、V, 6, 4 与 VI, 2, 10。

测度的虚部：III, 1, 5 与 VI, 2, 8。

诱导测度，在局部紧子空间上的：IV, 5, 7 与 VI, 2, 10。

诱导测度，在开子集上的：III, 2, 1。

归纳极限，局部凸空间归纳系统的：见直接极限。

Hölder 不等式：I, 2 与 IV, 6, 4。

Hardy 不等式：IV, 6，习题 19。

M. Riesz 不等式：IV, 6，习题 18。

Minkowski 不等式：I, 2 与 IV, 3, 1。

平均不等式：IV, 6, 2。

下确界：II, 1, 1，脚注。

内测度：IV, 4，习题 7。

可积函数：IV, 4, 1。

可积函数，\(p\) 次幂：IV, 3, 4。

可积分，\(\mu\) -可积分：IV, 4, 5。

函数在 A 上的积分（或延拓到 A 上的积分）：V, 5, 3。

标量本质可积映射的积分：V, 3, 1、VI, 1, 1 与 VI, 2, 10。

本质可积函数的积分：V, 1, 3。

积分，二重：III, 4, 1。

积分，本质上：V, 1, 1。

积分，下：IV, 4，习题 5。

积分，多重、n 重：III, 4, 4。

积分，具紧支集的连续数值函数的：III, 1, 3。

积分，数值函数关于向量测度的：VI, 2, 2。

积分，向量值函数关于正测度的：IV, 4, 1、V, 1, 3 与 VI, 1, 1。

积分，向量值函数关于向量测度的：VI, 2, 7。

积分，标量意义具紧支集的弱连续函数的：III, 3, 1。

积分，可积函数的：IV, 4, 1。

积分，上：IV, 1, 1、IV, 1, 3 与 IV, 4，习题 5。

分部积分：V, 8，习题 9。

不变测度，同胚下的：III, 1, 3。

局部同胚下测度的逆像：V, 6, 6。

逆极限，测度逆系统的：III, 4, 5。

逆系统（测度的）：III, 4, 5。

孤立子空间，Riesz 空间的：II, 1，习题 4。

Krein 定理：IV, 7，习题 10。

格序向量空间：II, 1, 1。

格线性映射：II, 2，习题 5。

Lebesgue 函数，第 \(n\) 个：IV, 6，习题 15。

Lebesgue 测度：III, 1, 3、III, 4, 1 与 III, 4, 4。

Lebesgue 分解定理：V, 5, 7。

Lebesgue 定理：IV, 3, 7 与 IV, 4, 3。

Lebesgue--Fubini 定理：V, 8, 4。

Lebesgue--Nikodym 定理：V, 5, 5。

提升性质：VI, 2, 5。

序列的极限分布：III, 2，习题 5。

上极限、下极限，完全格序空间中的：II, 1，习题 9。

极限，直接（或归纳）：III, 1, 1。

极限，逆（或射影）：III, 4, 5。

极限，完全格序空间中滤子的：II, 1，习题 9。

线性形式，正：II, 2, 1。

线性形式，相对有界：II, 2, 2。

局部化原理，可测函数的：IV, 5, 2。

局部化原理，测度的：III, 2, 1。

局部几乎处处：IV, 5, 2。

局部可数（子集的集）：IV, 5, 9。

局部可数的 \(\geqslant 0\) 函数族：V, 5, 4。

局部可积函数：V, 5, 1。

在 A 上局部可积的函数：V, 5, 3。

局部可忽略（\(\mu\) -可忽略）函数：IV, 5, 2。

局部可忽略集，局部 \(\mu\) -可忽略集：IV, 5, 2。

下积分：IV, 4，习题 5。

\(\mu\) -适应对：V, 4, 1。

\(\mu\) -适宜映射：V, 3, 1 与 VI, 1, 1，脚注。

\(\mu\) -稠密（紧子集的集）：IV, 5, 8。

\(\mu\) -等价函数：IV, 2, 4。

\(\mu\) -可积函数：IV, 4, 1。

\(\mu\) -可积分：IV, 4, 5。

\(\mu\) -最大值、\(\mu\) -最小值：IV, 6, 2。

\(\mu\) -可测（集、函数）：IV, 5, 1。

\(\mu\) -可测等价关系：VI, 3, 4。

\(\mu\) -适度（函数、子集）：V, 1, 2。

\(\mu\) -可忽略函数：IV, 2, 1 与 IV, 2, 4。

\(\mu\) -可忽略集：IV, 2, 2。

\(\mu\) -预适宜映射：V, 3, 1。

\(\mu\) -逆紧映射：V, 6, 1、V, 6, 4 与 VI, 2, 10。

可优超测度：VI, 2, 3。

映射，\(\mu\) -适宜：V, 3, 1。

映射，\(\mu\) -预适宜：V, 3, 1。

映射，\(\mu\) -逆紧（或对 \(\mu\) 逆紧）：V, 6, 1、V, 6, 4 与 VI, 2, 10。

映射，保守：V, 7，习题 9。

映射，标量意义具紧支集：III, 3, 1。

映射，弱连续：III, 3, 1。

总质量，有界测度的：III, 1, 8 与 IV, 4, 7。

测度，由质量定义的：III, 1, 3。

依测度最大值，\(\mu\) -最大值：IV, 6, 2。

函数的平均（值）：III, 1, 8。

\(p\) 阶平均收敛：IV, 3. 3。

平均收敛：IV, 3, 3。

平均不等式：IV, 6, 2。

可测（集、函数）：IV, 5, 1。

可测等价关系：VI, 3, 4。

定义在 X 的可测子集上的可测函数：IV, 5, 10 与 V, 5, 3。

在 A 上可测的函数：IV, 5, 10 与 V, 5, 3。

可测截面：VI, 3, 4。

可测阶梯函数：IV, 5, 5。

属于扩散定义域的测度：V, 3, 5。

集中于一集上的测度：V, 5, 7。

可积分的测度：IV, 4, 5。

以 \(\mu\) 为基的测度：V, 5, 2。

以密度 g（关于 \(\theta\) ）的测度：V, 5, 2。

测度，绝对值：III, 1, 6 与 VI, 2, 8。

测度，原子：III, 1, 3。

测度，有界：III, 1, 8 与 VI, 2, 9。

测度，复测度：III, 1, 3 与 VI, 2, 8。

测度，共轭：III, 1, 5 与 VI, 2, 8。

测度，由一点的单位质量定义的：III, 1, 3。

测度，由质量定义的：III, 1, 3。

测度，弥散：V, 5, 10。

测度，Dirac：III, 1, 3。

测度，离散：III, 1, 3。

测度，像：V, 6, 1、V, 6, 4 与 VI, 2, 10。

测度，诱导：III, 2, 1、IV, 5, 7 与 VI, 2, 10。

测度，内：IV, 4，习题 7。

测度，同胚下不变的：III, 1, 3。

测度，在开集中不含质量的：III, 2, 2。

测度，Lebesgue：III, 1, 3、III, 4, 1 与 III, 4, 4。

测度，可优超，q-可优超：VI, 2, 3。

测度，范数：III, 1, 8。

测度，外：IV, 1, 2 与 IV, 1, 4。

测度，点：III, 2, 4。

测度，正：III, 1, 5。

测度，乘积：III, 4, 1、III, 4, 4、III, 4, 6 与 VI, 2, 10。

测度，伪像：VI, 3, 2。

测度，按等价关系取的商：VI, 3, 5。

测度，实：III, 1, 5 与 VI, 2, 1。

测度，实部与虚部：III, 1, 5 与 VI, 2, 8。

测度，标量：VI, 2, 1。

测度，支集：III, 2, 2。

测度，向量：VI, 2, 1。

测度，以密度 \(g\) 关于 \(\mu\) 的：III, 1, 4、V, 5, 2 与 VI, 2, 8。

测度，互奇异的：V, 5, 7。

测度，等价的：V, 5, 6 与 VI, 2, 8。

依测度最小值，\(\mu\) -最小值：IV, 6, 2。

Minkowski 不等式：I, 2 与 IV, 3, 1。

适度函数，\(\mu-\)：V, 1, 2。

适度测度：V, 1, 2。

适度子集，\(\mu-\)：V, 1, 2。

多重积分：III, 4, 4。

\(n\) 重积分：III, 4, 4。

可忽略函数：IV, 2, 1 与 IV, 2, 4。

可忽略集：IV, 2, 2。

扩散的范数：V, 3, 5。

测度的范数：III, 1, 8。

数值函数：I, 1, 1，脚注。

\(n\) 阶多重积分：III, 4, 4。

p 阶平均收敛：IV, 3, 3。

p 阶等度可积（函数集）：IV, 5, 11。

序直和：II, 1, 4。

标准正交序列，\(\mathcal{L}^2\) 中的：IV, 6，习题 15。

标准正交系，Haar：IV, 6，习题 17。

外测度：IV, 1, 2 与 IV, 1, 4。

\(p\) 次幂可积函数：IV, 3, 4。

测度的实部与虚部：III, 1, 5 与 VI, 2, 8。

\(\Phi\) -阶梯函数：IV, 4, 9。

点测度：III, 2, 4。

点，\(\mathcal{H}\) -极值：IV, 7, 3。

正线性形式：II, 2, 1。

正测度：III, 1, 5。

预适宜映射，\(\mu-\)：V, 3, 1。

可测函数的局部化原理：IV, 5, 2。

测度的局部化原理：III, 2, 1。

正测度族的乘积：III, 4, 6。

有限个测度的乘积：III, 4, 4。

测度与连续函数的乘积：III, 1, 4。

测度与局部可积函数的乘积：V, 5, 2。

两测度的乘积：III, 4, 1 与 VI, 2, 10。

射影极限，测度射影系统的：见逆极限。

测度的射影系统：见逆系统。

逆紧映射，\(\mu-\)：V, 6, 1、V, 6, 4 与 VI, 2, 10。

性质（GDF）：VI, 1, 4。

几乎处处成立的性质：IV, 2, 3。

局部几乎处处成立的性质：IV, 5, 2。

性质，提升：VI, 2, 5。

伪像（类、测度）：VI, 3, 2。

\(q\) -可优超测度：VI, 2, 3。

可求积集：IV, 5，习题 17。

均方收敛：IV, 3, 3。

拟可积函数：IV, 5，习题 15。

拟强拓扑：III, 1，习题 8。

商测度，测度按等价关系的：VI, 3, 5。

实测度：III, 1, 5 与 VI, 2, 1。

测度的实部：III, 1, 5 与 VI, 2, 8。

重排，函数的递减：IV, 5，习题 29。

沿滤子退向无穷：III, 2, 2。

相对有界线性形式：II, 2, 2。

测度在局部紧子集上的限制：IV, 5, 7 与 VI, 2, 10。

测度在开集上的限制：III, 2, 1。

Riesz 空间：II, 1, 1。

F. Riesz 定理：II, 1, 5。

M. Riesz 不等式：IV, 6，习题 18。

标量测度：VI, 2, 1。

标量本质可积映射：V, 3, 1、VI, 1, 1 与 VI, 2, 10。

标量意义以 \(\mu\) 为基，向量测度：VI, 2, 5。

标量意义具紧支集，映射：III, 3, 1。

标量良可积函数：VI, 1，习题 19。

标量意义，函数具有某性质：VI, 1, 1。

截面，可测：VI, 3, 4。

可分：IV, 2, 4，脚注。

几乎处处收敛的序列：IV, 2, 5。

关于一测度等分布的序列：III, 2，习题 5。

几乎处处收敛的级数：IV, 2, 5。

集函数，加性：IV, 4, 9。

切片，测度按之的分解：V, 6, 6 与 VI, 3, 1。

空间，完全格序：II, 1, 3。

空间，Riesz：II, 1, 1。

阶梯函数：IV, 4, 9。

阶梯函数，\(\Phi-\)：IV, 4, 9。

阶梯函数，可测：IV, 5, 5。

Stieltjes 测度：V, 6，习题 5。

Stone 空间：II, 1，习题 13。

严格紧收敛（在测度空间中）：III, 1, 10。

严格紧集，\(\mathcal{K}(\mathrm{X};\mathrm{E})\) 中的：III, 1, 1。

正测度的可和族：V, 2, 1。

函数的支集：III, 1, 1。

测度的支集：III, 2, 2。

向量测度的支集：VI, 2, 1。

上确界：II, 1, 1，脚注。

Bishop 定理：IV, 7, 5。

可数凸性定理：IV, 3, 2。

Egoroff 定理：IV, 5, 4。

F. Riesz 定理：II, 1, 5。

Krein 定理：IV, 7，习题 10。

Lebesgue 定理：IV, 3, 7 与 IV, 4, 3。

定理，Dunford--Pettis：VI, 2, 5。

定理，Gelfand--Dunford：VI, 1, 4。

定理，Lebesgue 分解：V, 5, 7。

定理，Lebesgue--Fubini：V, 8, 4。

定理，Lebesgue--Nikodym：V, 5, 5。

Choquet 定理：IV, 7, 2 与 IV, 7, 6。

紧收敛拓扑、严格紧收敛拓扑（在测度空间上）：III, 1, 10。

平均收敛、\(p\) 阶平均收敛、均方收敛的拓扑：IV, 3. 3。

依测度收敛的拓扑：IV, 5, 11。

拓扑，拟强：III, 1，习题 8。

拓扑，超强：III, 1，习题 15。

拓扑，淡：III, 1, 9。

总质量，有界测度的：III, 1, 8 与 IV, 4, 7。

子集的 σ-集环（σ-代数）：IV, 4, 9，脚注。

几乎处处成立，性质：IV, 2, 3。

局部几乎处处成立，性质：IV, 5, 2。

超强拓扑：III, 1，习题 15。

依测度收敛的一致结构：IV, 5, 11。

普遍可测（集、函数）：V, 3, 4。

上积分：IV, 1, 1、IV, 1, 3 与 IV, 4，习题 5。

上积分，本质：V, 1, 1。

淡收敛：III, 1, 9。

淡拓扑：III, 1, 9。

淡 \(\mu\) -可测函数：V, 3, 1。

淡连续函数：V, 3, 1。

向量测度：VI, 2, 1。

以 \(\mu\) 为基的向量测度：VI, 2, 4。

向量测度，\(q\) -可优超：VI, 2, 3。

向量测度，可优超：VI, 2, 3。

向量测度，标量意义以 \(\mu\) 为基：VI, 2, 5。

向量测度，支集：VI, 2, 1。

弱连续映射：III, 3, 1。

良可积，标量：VI, 1，习题 19。

